% ======================================================================================
% Nom     : reports/latex/sections/06_algorithmes.tex
% Rôle    : Section présentant les algorithmes implémentés.
% Auteur  : Maxime BRONNY
% Version : V1
% Cadre   : Réalisé dans le cadre du cours Techniques d'apprentissage artificiel
%           Master 1 Informatique Big Data
% Usage   : Fichier inclus dans main.tex lors de la compilation du rapport.
% ======================================================================================

\section{Algorithmes implémentés}
\label{sec:algorithmes}

Les trois algorithmes sont implémentés from scratch en Python, avec NumPy pour
le calcul vectoriel, dans le module \texttt{src/models/}. Nous présentons pour
chacun le principe, les équations, les choix d'implémentation et les limites.

\subsection{Régression logistique}
\label{subsec:lr}

\paragraph{Principe.}
La régression logistique \cite{hosmer2000logistic} modélise directement la
probabilité de défaut comme une fonction du score linéaire
$z = \mathbf{w} \cdot \mathbf{x} + b$ passée dans la sigmoïde :
\[
  p(y = 1 \mid \mathbf{x}) = \sigma(z) = \frac{1}{1 + e^{-z}} .
\]
Les paramètres $\mathbf{w}$ et $b$ sont appris en minimisant l'entropie croisée
(\emph{log-loss}) sur l'ensemble d'entraînement :
\[
  J(\mathbf{w}, b) = -\frac{1}{n} \sum_{i=1}^{n}
    \left[ y_i \ln p_i + (1 - y_i) \ln (1 - p_i) \right],
\]
fonction convexe dont le minimum global est atteignable par descente de
gradient. Les gradients ont une forme remarquablement simple :
\[
  \frac{\partial J}{\partial \mathbf{w}} = \frac{1}{n} X^{\top}(\mathbf{p} - \mathbf{y}),
  \qquad
  \frac{\partial J}{\partial b} = \frac{1}{n} \sum_{i=1}^{n} (p_i - y_i).
\]

\paragraph{Implémentation.}
Notre implémentation (\texttt{src/models/logistic\_regression.py}) utilise une
descente de gradient \emph{full batch} : à chaque itération, le gradient est
calculé sur tout l'ensemble d'entraînement en une seule opération matricielle.
Deux détails numériques méritent mention : la sigmoïde est calculée par deux
formules équivalentes selon le signe de $z$ afin d'éviter tout dépassement de
capacité de l'exponentielle, et les probabilités sont bornées avant le
logarithme pour éviter $\ln 0$. Les hyperparamètres d'optimisation sont un taux
d'apprentissage de $0{,}1$ et $1\,000$ itérations, sur données normalisées. La
figure~\ref{fig:lr-cost} montre la décroissance régulière de la log-loss : la
descente de gradient converge sans oscillation, signe d'un taux d'apprentissage
bien dimensionné.

\begin{figure}[H]
  \centering
  \includegraphics[width=0.68\textwidth]{lr_cost_curve.png}
  \caption{Convergence de la descente de gradient : décroissance de la
    log-loss sur l'ensemble d'entraînement ($1\,000$ itérations, taux
    d'apprentissage $0{,}1$).}
  \label{fig:lr-cost}
\end{figure}

Le seuil qui convertit la probabilité en classe est traité comme un
hyperparamètre choisi sur la validation ; la valeur retenue, $0{,}3$, est
inférieure au seuil naïf de $0{,}5$ : en présence de classes déséquilibrées,
abaisser le seuil augmente le rappel, ce qui est cohérent avec le coût élevé
des faux négatifs.

\paragraph{Limites.}
La frontière de décision reste un hyperplan : les interactions entre variables
(un même taux d'endettement n'a pas le même sens selon le grade du dossier) ne
sont pas modélisées. Le modèle est de plus pénalisé par l'encodage ordinal des
variables nominales, qui lui impose un ordre arbitraire.

\subsection{Arbre de décision CART}
\label{subsec:dt}

\paragraph{Principe.}
L'algorithme CART \cite{breiman1984cart} construit un arbre binaire par
partitionnement récursif. À chaque n\oe ud, il cherche le couple (variable,
seuil) qui sépare le mieux les exemples, au sens de l'impureté de Gini : pour
un groupe contenant une proportion $p$ de défauts, $G = 2p(1-p)$, nulle pour un
groupe pur et maximale pour un mélange $50/50$. Le couple retenu minimise la
moyenne pondérée des impuretés des deux groupes enfants. La construction
s'arrête quand le n\oe ud est pur, trop petit, ou que la profondeur maximale
est atteinte ; chaque feuille prédit alors la proportion de défauts des
exemples qu'elle contient, interprétable comme une probabilité.

\paragraph{Implémentation.}
La difficulté pratique est le coût de la recherche du meilleur seuil. Notre
implémentation (\texttt{src/models/decision\_tree.py}) la vectorise : pour
chaque variable, les valeurs sont triées et les sommes cumulées des étiquettes
donnent, pour \emph{toutes} les positions de coupure possibles à la fois, la
proportion de défauts de part et d'autre, et donc le Gini de tous les seuils
candidats en une seule passe NumPy, sans boucle Python sur les seuils. Cette
vectorisation rend l'entraînement quasi instantané sur 22\,802 exemples. Deux
garde-fous limitent le surapprentissage : un n\oe ud n'est divisé que s'il
contient au moins 20 exemples, et toute feuille doit en contenir au moins 5.
La profondeur maximale, hyperparamètre principal, est choisie sur la
validation ; l'arbre retenu a une profondeur de $7$ et compte $105$ n\oe uds.

\paragraph{Limites.}
Un arbre seul est un modèle instable : une perturbation des données
d'entraînement peut changer les premières divisions et, par effet de cascade,
toute la structure. Sans limite de profondeur, il mémoriserait le bruit du
train ; le contrôle par la validation est donc indispensable. Enfin, sa
frontière de décision « en escalier » approxime grossièrement les relations
continues douces.

\subsection{$k$ plus proches voisins ($k$-NN)}
\label{subsec:knn}

\paragraph{Principe.}
Le $k$-NN \cite{cover1967knn} ne construit aucun modèle : l'« apprentissage »
consiste à mémoriser l'ensemble d'entraînement. Pour classer un nouvel exemple,
on calcule sa distance euclidienne à tous les exemples mémorisés, on retient
les $k$ plus proches et l'on prédit la proportion de défauts parmi eux. Le
paramètre $k$ règle le compromis biais-variance : petit, la décision est locale
mais bruitée ; grand, elle est lissée mais peut noyer les structures locales.

\paragraph{Implémentation.}
La difficulté est ici la mémoire : la matrice complète des distances entre les
4\,886 exemples de test et les 22\,802 exemples d'entraînement occuperait de
l'ordre de 900~Mo. Notre implémentation (\texttt{src/models/knn.py}) procède
donc par paquets de quelques centaines d'exemples. Deux optimisations
classiques sont utilisées : le développement
$\lVert a - b \rVert^2 = \lVert a \rVert^2 + \lVert b \rVert^2 - 2\,a \cdot b$,
dont le terme $\lVert a \rVert^2$, constant pour un exemple donné, peut être
omis car il ne change pas l'ordre des voisins ; et la sélection des $k$ plus
petites distances par \texttt{np.argpartition}, en temps linéaire, au lieu d'un
tri complet. La valeur retenue sur la validation est $k = 31$.

\paragraph{Limites.}
Toute la charge de calcul est reportée à la prédiction, qui parcourt
l'ensemble d'entraînement complet pour chaque exemple, ce qui est rédhibitoire
pour un système temps réel à fort trafic sans structures d'index dédiées. La méthode
est par ailleurs très sensible à l'échelle des variables (d'où la
normalisation) et sa pertinence se dégrade en haute dimension.
